\section{A diagonal estimate for the Hodge rank}\label{sec:diagonal}

We now show that a lower bound for local positivity of a normalized
polarization bounds $r(U)$.  The geometric part of the proof constructs a
K\"ahler class on a modification of the blown-up diagonal, with all
uncontrolled corrections supported on subvarieties that fail to dominate
one factor.  Lemma~\ref{hodge:basic} makes these corrections invisible to
$T(U)$.  The resulting diagonal class has a fixed nonzero intersection
but a norm that decreases with $r(U)$; Lemma~\ref{lem:hodge-gap} converts
this comparison into the required bound.

For an ample rational divisor $h$ on a projective variety and a smooth
point $x$, its Seshadri constant is
\[
 \epsilon(h;x)=\inf_{C\ni x}\frac{h\cdot C}{\mult_x C},
\]
where the infimum runs over integral curves through $x$.  Equivalently,
it is the largest $s$ for which $\sigma^*h-sE_x$ is nef on the blowup
$\sigma$ of $x$.

\begin{lemma}[Diagonal estimate]\label{lem:diagonal}
Fix an integer $n\ge2$ and a real number $c>0$.  Let $X$ be a normal
projective complex $n$-fold with canonical singularities and
$K_X\sim0$ as an integral principal divisor.  Suppose $h$ is an ample
rational divisor such that
\[
 1\le h^n\le2,
 \qquad \epsilon(h;x)>c
 \quad\text{at very general smooth points }x\in X.
\]
Then $r(U)=\dim_\Q T(U)$, for any smooth projective resolution $U$ of
$X$, is bounded by a constant depending only on $n$ and $c$.
\end{lemma}

\begin{proof}
Choose a resolution $\rho:U\to X$, isomorphic over the smooth locus,
obtained by blowups, and an effective exceptional divisor $N$ such that
$-N$ is $\rho$-ample.  The canonical form trivializing $K_X$ gives
\[
 K_U\ge0,\qquad K_U\text{ exceptional over }X.
\]
Put $H=\rho^*h$ and $l=H^n\in[1,2]$.  We also denote by $H$ a smooth
closed semipositive representative, obtained by pulling back a
Fubini--Study form from an embedding given by a multiple of $h$.
Let
\[
 Y=U\times U,\qquad
 \pi:P=\operatorname{Bl}_{\Delta}Y\longrightarrow Y
\]
be the blowup of the diagonal, with exceptional divisor $E$.  Subscripts
$1,2$ indicate pullback from the two factors, on every subsequent model.
A subvariety of any such model will be called \emph{double dominant} if
its maps to both copies of $U$ are dominant.

\subsection*{Diagonal conditions and the number of sections}

Fix rational numbers
\[
 t=\frac14,\qquad \tau=n-1+t,\qquad L=10n^2,
\]
and choose a rational $a_2$ depending only on $n,c$, sufficiently large
that $a_2>L/c$ and $a_2>\tau(1+1/(10n))$.  Set $a_1=10na_2$ and
$A=a_1H_1+a_2H_2$.  These choices give the strict inequalities
\begin{equation}\label{diag:parameter-margins}
 a_2>\tau(1+a_2/a_1),
 \qquad
 \frac{n}{1+n/L}>\tau(1+a_2/a_1).
\end{equation}
They remain true after sufficiently small rational perturbations of
$a_1,a_2,t$.  We first prove that
\begin{equation}\label{diag:base-locus-goal}
 G=A-\tau E\text{ is big},
 \qquad
 \mathbf B(A-tE)\text{ has no double-dominant component},
\end{equation}
where $\mathbf B$ denotes the stable base locus.

For sufficiently divisible positive integers $m$, put
\[
 Q_m=H^0(Y,mA),\qquad
 V_m=H^0(P,mG)
     =H^0(Y,mA\otimes\mathcal I_\Delta^{m\tau})\subset Q_m.
\]
The conormal bundle of the diagonal is $\Omega_U^1$.  Its infinitesimal
layers therefore give
\begin{equation}\label{diag:layers}
 \dim(Q_m/V_m)
 \le\sum_{0\le j<m\tau}
 h^0\bigl(U,\OO_U(m(a_1+a_2)H)\otimes\Sym^j\Omega_U^1\bigr).
\end{equation}
We need an estimate uniform in $j$, because $j$ grows with $m$.

The reflexive cotangent sheaf of $X$ is slope semistable of degree zero
with respect to $h$.  This follows from generic semipositivity for the
cotangent bundle and $K_X\equiv0$; the precise singular-variety
semistability and restriction statements we use are
\citep[Theorem~5.3 and Proposition~5.4]{GKP2016}; see also
\citep[Theorem~1.3]{CP2019}.  Choose a fixed sufficiently positive
general complete-intersection flag on $U$, cut by pullbacks of
hypersurfaces on $X$, terminating in a smooth curve $C$ avoiding the
exceptional locus.  Its members are smooth and integral.  The bundle
$\Omega_U^1|_C$ is semistable of degree zero, as is every symmetric
power in characteristic zero.

Write $D_0=m_0H$, with $m_0h$ very ample, and
$r_j=\binom{j+n-1}{n-1}$.  On $C$, a negative twist of
$\Sym^j\Omega_U^1|_C$ has no sections.  For $k\ge0$, subtracting
$k\deg_C D_0+1$ points gives
\[
 h^0\bigl(C,\OO_C(kD_0)\otimes\Sym^j\Omega_U^1|_C\bigr)
 \le r_j(k\deg_C D_0+1).
\]
Negative twists also have no sections on higher flag members: general
curves of this kind sweep out a dense subset, and their restrictions
have negative slope.  If a hypersurface cut in the flag has degree $e$
with respect to $D_0$, write $b_{Z,j}(k)$ for the section dimension on
the member $Z$.  The restriction sequences and negative-twist vanishing give
\[
 b_{Z,j}(k)\le
 \sum_{a=0}^{\lfloor k/e\rfloor}b_{Z',j}(k-ae),
 \qquad Z'\in|eD_0|_Z.
\]
Since $D_0^{i-1}\cdot Z'=e(D_0^i\cdot Z)$, summing the degree-$(i-1)$
leading term on $Z'$ divides its coefficient by $ei$.  Induction gives,
on an $i$-dimensional flag member $Z$,
\[
 h^0\bigl(Z,\OO_Z(kD_0)\otimes\Sym^j\Omega_U^1|_Z\bigr)
 \le r_j\left(\frac{(D_0^i\cdot Z)k^i}{i!}
                +O((k+1)^{i-1})\right),
\]
with the implied constant independent of $j$.  Since
$\sum_{j<q}r_j=q^n/n!+O(q^{n-1})$, \eqref{diag:layers} becomes
\begin{equation}\label{diag:section-upper}
 \dim(Q_m/V_m)
 \le \frac{l(a_1+a_2)^n\tau^n}{(n!)^2}m^{2n}
       +o(m^{2n}).
\end{equation}
All asymptotic errors in this proof concern the fixed variety and fixed
divisors; their thresholds need not be uniform in $X$.
On the other hand,
\[
 \dim Q_m=
 \frac{l^2a_1^na_2^n}{(n!)^2}m^{2n}+o(m^{2n}).
\]
The first inequality in \eqref{diag:parameter-margins}, together with
$l\ge1$, makes the difference of the leading coefficients positive.
Thus $G$ is big.

We next show that, for all sufficiently large divisible $m$, there is no
divisorial valuation $v$ over $Y$ whose centre is double dominant and
satisfies
\begin{equation}\label{diag:excluded-valuation}
 v(V_m)\ge m A_Y(v).
\end{equation}
Here $v(V_m)$ is the order of the base ideal of $V_m\subset H^0(Y,mA)$,
and $A_Y(v)$ is the log discrepancy.  The Seshadri assumption and
$a_2c>L$ imply that $ma_2H$ separates jets modulo the $Lm$-th power of
the maximal ideal on a dense open subset of $U$, for all sufficiently
large divisible $m$.  Indeed, at a smooth point satisfying the
Seshadri bound the divisor $a_2h-LE_x$ on its blowup is ample.  Serre
vanishing supplies the jet separation, and surjectivity of evaluation
is open.

If \eqref{diag:excluded-valuation} held, localize a resolution carrying
$v$ over the generic point of the first factor and extend to
$\overline{\C(U)}$.  The divisor remains horizontal; its discrepancy
and its base-ideal order become the corresponding quantities on this
smooth geometric generic fibre.  Double dominance ensures that its
centre meets the second-factor jet-separation open.  At a closed point
$y$ of that intersection, the restricted base ideal $I$ would satisfy
$\lct_y(I)\le1/m$.  For
$J=I+\mathfrak m_y^{Lm}$, the sum-of-ideals inequality gives
\[
 \lct_y(J)\le\lct_y(I)+\frac{n}{Lm}
             \le\frac{1+n/L}{m}
\]
\citep[Lemma~2.18]{dFEM2011}.  The threshold--colength inequality
\citep[Theorem~1.1]{dFEM2004} then gives
\begin{equation}\label{diag:colength}
 s:=\operatorname{length}(\OO_y/J)
 \ge\frac{m^n}{n!}\left(\frac{n}{1+n/L}\right)^n.
\end{equation}

Jet separation allows us to choose $s$ sections $g_1,\ldots,g_s$ of
$ma_2H$, defined over $\C$, whose images modulo $J$ are linearly
independent over $\overline{\C(U)}$.  Namely, after scalar extension the
images of any complex basis span the quotient, so one can select a
basis among them.  If $f_1,\ldots,f_q$ is a complex basis of
$H^0(U,ma_1H)$, their $qs$ products $f_i g_j$ are linearly independent
modulo $V_m$.  To see this, restrict a complex linear relation to the
geometric generic fibre and reduce modulo $J$.  Independence of the
$g_j$ there makes each coefficient $\sum_i c_{ij}f_i$ zero at the
generic point.  These are complex linear relations among the $f_i$,
so all $c_{ij}$ vanish.

Consequently \eqref{diag:colength} gives the lower leading coefficient
\[
 \frac{l a_1^n}{(n!)^2}
       \left(\frac{n}{1+n/L}\right)^n
\]
for $m^{-2n}\dim(Q_m/V_m)$.  By the second inequality in
\eqref{diag:parameter-margins}, this exceeds the upper coefficient in
\eqref{diag:section-upper}, a contradiction.  The numerical bounds are
uniform over the chosen jet-separation open, so the sufficiently large
$m$ can be fixed independently of the hypothetical valuation.

This exclusion controls the discrepancies of a general divisor with the
prescribed diagonal vanishing along every double-dominant centre.  We now
turn that discrepancy control into the second assertion of
\eqref{diag:base-locus-goal}, by comparing an adjoint ample model with the
diagonal blowup near such centres.

\subsection*{Removing base components that dominate both factors}

Fix such an $m$, also large enough that $|mG|$ has generically finite
map.  First choose a log resolution $p_0:W\to P$ of its base ideal,
together with $E$ and $N_K=\pi^*K_Y$.  Then choose
$\Delta_P\in\frac1m|mG|$ general.  Bertini makes its free part smooth
and transverse to the fixed normal crossings divisor, so the resulting
boundary on $W$ is log smooth.  Write
\[
 K_W+\Delta_W=p_0^*(K_P+\Delta_P).
\]
Every double-dominant divisor on $W$ has coefficient less than one in
$\Delta_W$.  To verify this, use
\begin{equation}\label{diag:canonical-blowup}
 K_P+\Delta_P
  =\pi^*(K_Y+\pi_*\Delta_P)-tE.
\end{equation}
For the finitely many exceptional or fixed divisors on this resolution,
a general member has valuation $v(V_m)/m$ after pushforward to $Y$.
Its coefficient in $\Delta_W$ is therefore
\[
 1-A_Y(v)+\frac{v(V_m)}m-t v(E)<1
\]
when its centre is double dominant, by the exclusion of
\eqref{diag:excluded-valuation}.  Components of the general free part
have coefficient $1/m<1$.  Formula \eqref{diag:canonical-blowup} remains
valid if $E$ occurs in the fixed part of $|mG|$: pulling back the
pushforward restores exactly the prescribed multiplicity $m\tau E$.

Replace negative coefficients of $\Delta_W$ by zero, and decrease all
coefficients at least one to rational numbers strictly below one.
The resulting boundary $B_W$ is klt.  Set $L_W=K_W+B_W$; then
\begin{equation}\label{diag:adjusted-boundary}
 K_W+B_W=p_0^*(D'+N_K)+G_e-F_N,
 \qquad D'=K_P+\Delta_P-N_K\sim_\Q A-tE,
\end{equation}
where $N_K\ge0$, $G_e\ge0$ is $p_0$-exceptional, and every
component of $F_N\ge0$ fails to be double dominant.  The assertion about
$G_e$ follows because $\Delta_P$ is effective: negative coefficients
can occur only on $p_0$-exceptional divisors.  The assertion about $F_N$
is the preceding coefficient bound.

The divisor $K_W+B_W$ is big.  Indeed, $K_W$ is effective, and $B_W$
contains with positive coefficient the big free part of $|mG|$.
The existence of ample models for big klt adjoints
\citep{BCHM2010} gives a projective klt pair $(V,B_V)$, the ample
rational adjoint divisor $D_V=K_V+B_V$, and a common smooth resolution
\[
 r_0:S\to W,\qquad q:S\to V,\qquad p=p_0r_0:S\to P
\]
such that
\begin{equation}\label{diag:ample-model}
 r_0^*(K_W+B_W)=q^*D_V+E_V,
 \qquad E_V\ge0\text{ is }q\text{-exceptional}.
\end{equation}
We choose divisor representatives compatibly throughout these formulas.

Suppose that $\mathbf B(A-tE)$ has a double-dominant component.  Its
image in $Y$ contains a point $z=(x,y)$ satisfying two conditions:
neither $x$ nor $y$ lies on the image of a nonconstant rational curve in
$U$, and $z$ avoids the images of $N_K$ and $F_N$.  Such a point exists.
The effective canonical divisor makes $U$ non-uniruled, so images of
rational curves are contained in a countable union of proper closed
subsets.  Double dominance, and uncountability of $\C$, allow us to
avoid their inverse images and the finitely many other indicated
subsets.

Every fibre of $q$ is rationally chain connected because $(V,B_V)$ is
klt, by
\citep[Corollary~1.6]{HM2007}.  If a fibre meets $(\pi p)^{-1}(z)$,
all its rational chains have constant image in both copies of $U$:
a nonconstant image would be a rational curve through $x$ or $y$.
Thus that entire fibre maps to $z$.  This uses the all-fibres statement;
chains through special points may equivalently be obtained by proper
specialization.

It follows that $V$ maps to $Y$ near the points lying over $z$.  Here
is the precise local argument.  Let $\Gamma\subset V\times Y$ be the
closure of the graph of the birational map $V\dashrightarrow Y$.
Every fibre of $\Gamma\to V$ meeting $V\times\{z\}$ is a singleton,
by the preceding paragraph and the map $S\to\Gamma$.  Properness makes
this projection finite near each such fibre, and finite birationality
over normal $V$ makes it an isomorphism there.  The image in $Y$ of its
non-isomorphism locus is closed and avoids $z$.  Removing it, and the
images of $N_K,F_N$, gives an open $Y^0\ni z$ for which the graph is an
open $V^0\subset V$, proper over $Y^0$, and
$S_{Y^0}=q^{-1}(V^0)$.

Over $Y^0$, Equations \eqref{diag:adjusted-boundary} and
\eqref{diag:ample-model} reduce to
\begin{equation}\label{diag:relative-models}
 p^*D'+r_0^*G_e=q^*D_V+E_V.
\end{equation}
The divisors $D'$ on $P_{Y^0}$ and $D_V$ on $V^0$ are relatively ample
over $Y^0$: the first has relative class $-tE$, and the second is the
restriction of an ample divisor.  Adding an effective exceptional
divisor to a birational pullback does not change its pushforward
section algebra, by normality.  The relative Proj of the two sides of
\eqref{diag:relative-models} therefore identifies
$P_{Y^0}\simeq V^0$.  On this common model, pushforward gives
$D'=D_V$, and \eqref{diag:relative-models} gives
$E_V=r_0^*G_e$.

For divisible $k$, the model formulas give the following map of complete
section spaces:
\begin{equation}\label{diag:section-transfer}
 \begin{aligned}
 H^0(V,kD_V)&\simeq H^0(W,kL_W)
   \xrightarrow{\ \cdot s_{kF_N}\ }H^0(W,k(L_W+F_N))\\
   &\simeq H^0(P,k(D'+N_K)).
 \end{aligned}
\end{equation}
The first isomorphism pulls sections to $S$, multiplies by the canonical
section of $kE_V$, and descends to $W$.  The last uses
$L_W+F_N=p_0^*(D'+N_K)+G_e$ and removes the effective exceptional
divisor $G_e$ by pushforward.  Large divisible multiples of $D_V$ are
globally generated.  Their images in \eqref{diag:section-transfer}
generate above $z$: there $F_N=N_K=0$, and
$E_V=r_0^*G_e$ shows that multiplication and removal of the exceptional
factors cancel.  Thus under $P_{Y^0}\simeq V^0$ the transported section
is exactly the original section of $kD_V$.
Finally $N_K$ is a fixed divisor and can be removed.  Indeed, under
$Y\to X\times X$, $K_Y$ is effective exceptional and $A$ is pulled
back from $X\times X$, so
\[
 H^0(Y,k(A+K_Y))=s_{kK_Y}H^0(Y,kA)
\]
for divisible $k$.  Its pullback $N_K$ has zero coefficient along $E$;
dividing by its canonical section therefore preserves the vanishing
condition imposed by $-ktE$.  Hence
\[
 H^0(P,k(A-tE+N_K))=s_{kN_K}H^0(P,k(A-tE)).
\]
Generation above $z$ contradicts the assumed base component.  This
proves \eqref{diag:base-locus-goal}.

\subsection*{An ample class with controlled corrections}

The preceding argument controls where base conditions occur, rather
than their number or multiplicity.  This is enough: on each correction,
Lemma~\ref{hodge:basic} annihilates the pullback of $T(U)$ from at least
one factor.

Choose an ample rational divisor on $P$ of the form
\[
 J=m'(H_1+H_2)-N_1-N_2-\lambda E,
\]
with $m'$ large and $\lambda>0$ small.  For sufficiently small rational
$\varepsilon>0$, the stable base locus of $A-tE-\varepsilon J$ still
has no double-dominant component.  In fact
\[
 A-tE-\varepsilon J
 =(a_1-\varepsilon m')H_1+(a_2-\varepsilon m')H_2
   -(t-\varepsilon\lambda)E+\varepsilon(N_1+N_2).
\]
The first three terms satisfy \eqref{diag:base-locus-goal} by the strict
margins in \eqref{diag:parameter-margins}; adding the last effective
term can introduce base components only in its non-double-dominant
support.

Take a divisible multiple whose base locus is the stable base locus,
and principalize its base ideal on a smooth $w:\widetilde P\to P$,
with normal crossings support also including $E$.  This can be done
without changing the complement of the base locus.  Write
\[
 w^*(A-tE-\varepsilon J)=M+F_{\mathrm{fix}},
\]
where $M$ is semiample and $F_{\mathrm{fix}}\ge0$ has only
non-double-dominant components.  Choose an effective $w$-exceptional
$T$ with $-T$ relatively ample.  For sufficiently small rational
$u>0$, $\varepsilon w^*J-uT$ is ample.  Consequently
\[
 D=w^*(A-tE)-F_{\mathrm{fix}}-uT+H_1+H_2
\]
is ample and has a K\"ahler representative with $D\ge H_1+H_2$.
If $F$ is the strict transform of $E$, then
\begin{equation}\label{diag:ample-class}
 D=(a_1+1)H_1+(a_2+1)H_2-tF-Z,
 \qquad
 Z=t(w^*E-F)+F_{\mathrm{fix}}+uT\ge0.
\end{equation}
Its components are smooth and fail to be double dominant.  All further
constants $C$ depend only on $n,c$, and may increase from line to line.
Volume monotonicity in \eqref{diag:ample-class} gives
\begin{equation}\label{diag:volume}
 \int_{\widetilde P}D^{2n}
 \le\binom{2n}{n}(a_1+1)^n(a_2+1)^n l^2\le C.
\end{equation}

For general $x\in U$, the fibre $\widetilde P_x$ over the first
projection is smooth and birational to $U$.  All the correction
components avoid the fibre above the diagonal point $(x,x)$: for each
one, at least one projection has proper image.  Thus the original
exceptional space there is unchanged and
\begin{equation}\label{diag:fibre-intersections}
 F_x\simeq\PP^{n-1},\quad
 \OO(F_x)|_{F_x}=\OO_{\PP^{n-1}}(-1),\quad
 D|_{F_x}=t\,c_1(\OO_{\PP^{n-1}}(1)),\quad
 D_x^n\ge l,
\end{equation}
where $D_x=D|_{\widetilde P_x}$.  The last inequality follows from
$D_x\ge H_2|_{\widetilde P_x}$.

\subsection*{The diagonal class and its norm}

We will construct a middle-degree class $v$ whose intersection with $D^n$
is fixed, but whose norm is small when $r(U)$ is large.  A test class with
the same intersection and the improved growth of
Lemma~\ref{lem:hodge-gap} will then force an upper bound for $r(U)$.

Put $r=r(U)>0$; positivity follows from the canonical generator, since
$H^{n,0}(U)$ is one-dimensional.  By Lemma~\ref{hodge:basic}, $T(U)$
is primitive.  Choose a Hodge-homogeneous basis $\alpha_1,\ldots,\alpha_r$
of $T(U)_\C$, orthonormal for its Hodge--Riemann Hermitian form, so that
\[
 \int_U\alpha_i\overline{\alpha_j}=c_i\delta_{ij},
 \qquad |c_i|=1.
\]
Put $\beta_i=\overline{\alpha_i}/c_i$ and, on $\widetilde P$, set
\[
 v=\frac lr\sum_{i=1}^r\alpha_i\otimes\beta_i\in H^{n,n}(\widetilde P),
\]
using pullback from $Y$.  The graded product rule and the dual-basis
normalization give
\begin{equation}\label{diag:self-pairing}
 (-1)^n\int_{\widetilde P}v\overline v=\frac{l^2}{r},
 \qquad H_1v=H_2v=0.
\end{equation}
On the diagonal, the class $v$ has integral $l$; hence its restriction
to $F$ is $H_1^n|_F$.  On every component of $Z$, one projection has
image of dimension less than $n$, so Lemma~\ref{hodge:basic} kills the
corresponding factor in $v$.  The projection formula and
\eqref{diag:ample-class} now give
\begin{equation}\label{diag:intersection}
 Dv=-t[F]H_1^n,
 \qquad
 \int_{\widetilde P}vD^n=-t^n l.
\end{equation}
For the second identity, integrate over the first factor and use
$\int_{F_x}D_x^{n-1}=t^{n-1}$ from
\eqref{diag:fibre-intersections}.

For $R\ge1$, consider $\omega=D+RH_1$.  We claim that
\begin{equation}\label{diag:exceptional-norm}
 \norm{[F]H_1^n}_\omega\le C R^{-n/2}.
\end{equation}
By Hodge-star duality, it suffices to pair with harmonic
$(n-1,n-1)$-forms $g$.  On the general fibre, Lefschetz decomposition
of the $(1,1)$-class $[F_x]$ bounds its squared Hodge norm by a
dimensional constant times
\[
 |F_x^2D_x^{n-2}|+
 \frac{(F_xD_x^{n-1})^2}{D_x^n}.
\]
This is bounded by \eqref{diag:fibre-intersections}, since the two
numerators are $t^{n-2}$ and $t^{2n-2}$.  Consequently
\[
 \left|\int_{\widetilde P_x}[F_x]g|_{\widetilde P_x}\right|
 \le C\norm{g|_{\widetilde P_x}}_{L^2(D_x)}.
\]
Here replacing a closed restricted form by its harmonic representative
can only decrease its norm.  Integrate against $H^n$ on $U$ and apply
Cauchy--Schwarz.  Pointwise restriction to the vertical tangent space
does not increase the norm, while
\[
 H_1^n\wedge\omega^n\le C R^{-n}\omega^{2n}.
\]
Fibre integration outside the measure-zero exceptional sets therefore
gives
\[
 \int_U\norm{g|_{\widetilde P_x}}_{L^2(D_x)}^2H^n(x)
 \le C R^{-n}\norm{g}_\omega^2.
\]
Since $\int_UH^n=l\le2$, this proves
\eqref{diag:exceptional-norm}.

Decompose $v=v_0+v_1$ into its primitive and nonprimitive parts for
$\omega$.  Lefschetz linear algebra has a dimension-only inverse bound
on the nonprimitive middle-degree part.  Equations
\eqref{diag:self-pairing}, \eqref{diag:intersection}, and
\eqref{diag:exceptional-norm} imply
\[
 \norm{v_1}_\omega\le C\norm{\omega v}_\omega
                  \le C R^{-n/2}.
\]
Hodge--Riemann positivity on $v_0$, and orthogonality of Lefschetz
summands for both the Hodge norm and the middle intersection pairing,
then give
\begin{equation}\label{diag:small-norm}
 \norm v_\omega\le C\bigl(r^{-1/2}+R^{-n/2}\bigr).
\end{equation}
Indeed, the signed self-pairing of $v_0$ is its squared norm, whereas
the absolute self-pairing of $v_1$ is bounded by its squared norm;
\eqref{diag:self-pairing} supplies the remaining term $l^2/r$.

To test the fixed intersection in \eqref{diag:intersection}, set
\[
 z'=D^n-c'H_2^n,
 \qquad
 c'=\frac{\int H_1^nD^n}{\int H_1^nH_2^n}
    =\frac{\int H_1^nD^n}{l^2}.
\]
Thus $H_1^n z'=0$, and $|c'|\le C$ by $H_i\le D$ and
\eqref{diag:volume}.  The same inequalities give
$\norm{z'}_{D+H_1}\le C$: the indicated smooth representatives have
bounded pointwise norm for $D+H_1$, whose volume is bounded.  The form
$H_1$ has rank at most $n$, so Lemma~\ref{lem:hodge-gap} yields
\[
 \norm{z'}_{D+RH_1}\le C R^{(n-\delta_n)/2}.
\]
Since $vH_2=0$, cup-product Cauchy--Schwarz and
\eqref{diag:small-norm} now imply
\[
 t^n\le t^nl=\left|\int vz'\right|
 \le C\bigl(r^{-1/2}+R^{-n/2}\bigr)R^{(n-\delta_n)/2}.
\]
Choose $R$ depending only on $n,c$ so large that the term
$C R^{-\delta_n/2}$ is at most $t^n/2$.  The remaining inequality bounds
$r$ in terms of $n,c$, as required.
\end{proof}
